\begin{claim}{C.60}{Proved}{All positive homotopy groups of the relative one-partition kernel vanish}
Use the original plane loop, its generated group $G$,
the commuting operator group $H$, and the full global
quotient families of C.53.
Let $W$ be any $H$-invariant union of open $G$-orbits,
each having exactly one partition in that family.
Define $\mathcal K_W$ to be the subgroup of
$C_{\operatorname{Homeo}(\R^2)}(H)$ consisting of maps
which fix $\R^2\setminus W$ pointwise, preserve each
of those open orbits, and satisfy
\[
                 \Phi_z a|_O=\Phi_z
 \quad\hbox{for every actual global quotient of }O\subset W.
                                                               \tag{C.60.1}
\]
Give it the induced compact-open group topology.
Then
\[
                 \pi_n(\mathcal K_W,1)=0
                         \qquad(n\geq1).               \tag{C.60.2}
\]
More precisely, every based loop $a_t$ in this group
has a jointly continuous based contraction $R_{r,t}$
in the same group, with jointly continuous inverse,
and with the pointwise control
\[
 \begin{split}
 R_{r,t}(x)&\in\{a_s(x):0\leq s\leq t\},\\
 R_{r,t}^{-1}(x)&\in\{a_s^{-1}(x):0\leq s\leq t\}
 \qquad(0\leq r,t\leq1).
 \end{split}                                             \tag{C.60.3}
\]
Here $R_{0,t}=1$, $R_{1,t}=a_t$, and
$R_{r,0}=R_{r,1}=1$.
Every based sphere family has a based contraction
with the analogous control by its entire compact
family of point images and inverse images.
No compactness of any local quotient, finiteness of
operator-orbit representatives, or separate ambient-null
hypothesis is required.
\end{claim}

\begin{proof}
\emph{The exact winding input and the source covers.}
For $W=\varnothing$ the group is trivial.
Otherwise choose a global quotient of one of its
orbits. Its target is a nontrivial nonopen $G$-orbit,
so it lies outside $W$ and contains two distinct
points. Every $a_t$ fixes these two points.
The difference-direction map on
$\operatorname{Homeo}^+(\R^2)$ induces an isomorphism
on fundamental groups, as proved in C.32.7.
Its value on this loop is constant. Thus $a_t$ is
automatically null-homotopic in the ambient plane
homeomorphism group.
C.32 consequently makes its restriction to every
preserved open orbit null-homotopic in
$\operatorname{Homeo}(O)$.

Retain the original source universal covering
$p:X\to O$ and one original target covering $q$,
with $qF=\Phi p$ as in C.43--C.44.
The based lift $\widetilde a_t$ of $a_t|_O$ is a
closed loop of homeomorphisms of $X$:
the restricted null-homotopy just proved lifts,
so its endpoint is the identity deck map.
Equation (C.60.1) gives
$qF\widetilde a_t=qF$.
For a circle target the difference between the real
lifts is a continuous integer deck translation,
initially zero; for a line it is zero directly.
In both cases
\[
                         F\widetilde a_t=F.             \tag{C.60.4}
\]
This retains both original covers and their labels.
The construction below uses only this closed based
source lift and (C.60.4) as its winding input.

\emph{The actual affine fiber coordinates and operator equations.}
Put $Q_O=\overline{\operatorname{im}
(H_O\to\operatorname{Homeo}(O))}$, with its actual
topology and the restriction kernel of C.30.
The loop $a_t|_O$ commutes with $Q_O$.
Its based lift commutes with every element of the
all-lifts group $\widetilde Q_O$ of C.48.2:
the commutator is deck-valued, continuous in $t$,
and initially the identity.

Apply the established C.48 normal forms to this
proper group and to the unique actual partition.
The transitive base group used there is still the
one supplied by the original $G$-action; it has
not been replaced by the kernel loop $a_t$.
The resulting exact product chart $X\cong\R^2$
has $F(u,v)=v$, and every lifted operator is affine
in its fiber coordinate:
\[
 \widetilde q(u,v)=
   \bigl(\alpha_{\widetilde q}(v)u+
         \beta_{\widetilde q}(v),\,\overline q(v)\bigr),
                  \qquad \alpha_{\widetilde q}(v)\ne0. \tag{C.60.5}
\]
This displayed form includes all the cases
C.48.7--C.48.17: real and cyclic fiber kernels,
discrete and real base images, the affine
two-dimensional group, and the normalized lattice.
In particular it includes every source deck map.
No affine assertion is made for the original
ambient coordinates.

By (C.60.4) write
$\widetilde a_t(u,v)=(h_t(u,v),v)$.
For each $t,v$ the function $h_t(\,\cdot\,,v)$ is
an increasing onto line homeomorphism, and
$h_0(u,v)=h_1(u,v)=u$.
Define
\[
 \widetilde R_{r,t}(u,v)
                   =\bigl((1-r)u+r h_t(u,v),\,v\bigr). \tag{C.60.6}
\]
These are homeomorphisms with jointly continuous
inverses by strict monotonicity, properness, and
the two-point inverse bracketing argument in C.48.
Their operator equations are exact. Namely
commutation of $\widetilde a_t$ with (C.60.5) reads
\[
 h_t(\alpha(v)u+\beta(v),\overline q(v))
                         =\alpha(v)h_t(u,v)+\beta(v).
 \]
Taking the displayed convex combination with the
identity proves the same equation for
$\widetilde R_{r,t}$, including negative $\alpha$.
It commutes with the original deck group and
descends to $R^O_{r,t}$ on $O$.
Every other global real quotient is an exact
line reparametrization of $F$, by the dependent
case of C.45. It too is fixed, and descent through
its specified target cover proves (C.60.1).

Choose the chart on one representative of each
$H$-orbit of open orbits and put
$R^{hO}_{r,t}=hR^O_{r,t}h^{-1}$.
The representative construction commutes with
$Q_O$, hence with the actual image of $H_O$.
This makes transport independent of the operator
chosen, including its restriction kernel.
Set $R_{r,t}=1$ outside $W$.

\emph{Actual prefix tracks control all frontier limits.}
Fix $(u,v)$ and $t$. The real path
$s\mapsto h_s(u,v)$, $0\leq s\leq t$, begins at
$u$ and ends at $h_t(u,v)$.
Its image contains the whole closed interval
between these endpoints.
The first coordinate in (C.60.6) lies in that interval.
After projection this proves the first inclusion
in (C.60.3) on $O$.

For the inverse, write
$g(u)=(1-r)u+r h_t(u,v)$ with $v$ fixed.
Its strict monotonicity shows that $g^{-1}(u)$
lies between $u$ and $h_t^{-1}(u,v)$:
the values of $g$ at these two endpoints bracket $u$.
The continuous inverse prefix
$s\mapsto h_s^{-1}(u,v)$ contains that interval.
Projection proves the second inclusion in (C.60.3).
Operator transport preserves both statements
because the original loop commutes with $H$.
No continuous selection of the prefix parameter
$s$ is needed.

For any sequence $x_n\to x\notin W$, joint
continuity of the original loop and compactness
of its parameter interval give
\[
 \sup_{0\leq s\leq1}|a_s(x_n)-x|\longrightarrow0,\qquad
 \sup_{0\leq s\leq1}|a_s^{-1}(x_n)-x|\longrightarrow0.
                                                               \tag{C.60.7}
\]
Here $a_s(x)=a_s^{-1}(x)=x$ throughout.
The two track inclusions therefore prove joint
continuity of the constructed maps and inverses
at $x$, uniformly in $r,t$, regardless of which
open orbits contain $x_n$.
Inside $W$ a point has its open-orbit neighborhood,
where the intrinsic construction is already
jointly continuous.
Thus both global families are jointly continuous.
They are inverse orbitwise bijections and so
are homeomorphisms of the plane.
They commute with every original operator,
fix the exterior, preserve each orbit and every
global quotient, and satisfy all based-edge
equalities. This proves (C.60.2) for $n=1$.

\emph{Based sphere families in every higher degree.}
Let $n\geq2$ and let
$a:S^n\to\mathcal K_W$ be continuous, with
$a_{\lambda_*}=1$.
For each original source covering $p:X\to O$ the map
\[
 X\times S^n\longrightarrow O,\qquad
                  (u,\lambda)\longmapsto a_\lambda(pu)
                                                               \tag{C.60.8}
\]
has simply connected, locally path-connected domain.
The covering lifting criterion gives its unique
continuous lift with value $u_*$ at $(u_*,\lambda_*)$.
On $X\times\{\lambda_*\}$ that lift is the identity.
Lift the inverse family from the same basepoint.
Composing the two lifts gives a lift of $p$ on the
connected product and agrees with the identity at
the basepoint; uniqueness proves both inverse
identities. Each lifted map is therefore a
homeomorphism and both families are jointly continuous.

Commutation with every lifted operator and equality
of the real quotient coordinate follow as before:
the corresponding deck elements are continuous in
the connected parameter space and trivial at
$\lambda_*$. Apply (C.60.6) with $h_\lambda$ in place
of $h_t$. It is a based sphere contraction, preserves
all the exact operator and quotient equations, and
transports over $H$ just as above.

For any $\lambda$, choose a path in $S^n$ from
$\lambda_*$ to $\lambda$. Its lifted fiber evaluation
contains the interval between $u$ and
$h_\lambda(u,v)$. The forward interpolation and
inverse bracketing consequently give
\[
 \begin{split}
 R_{r,\lambda}(x)&\in
                       \{a_\mu(x):\mu\in S^n\},\\
 R_{r,\lambda}^{-1}(x)&\in
                       \{a_\mu^{-1}(x):\mu\in S^n\}.
 \end{split}                                             \tag{C.60.9}
\]
No continuous choice of those parameter paths is
used. Compactness of the whole sphere family
replaces the time interval in (C.60.7), proving
joint frontier continuity of maps and inverses
uniformly in $r,\lambda$.
The resulting homotopy is based at $\lambda_*$
and constant at $r=0$, so it descends to a
continuous filling on the cone $CS^n\cong D^{n+1}$.
This proves (C.60.2) in all positive degrees.

The conclusion supplies based contractions of every
sphere family in the relative one-partition quotient kernel,
including arbitrary accumulation and noncompact
local quotients. It does not supply a lift of a
nontrivial deformation of the quotient actions
or a decomposition of the original loop into
such a lift and this kernel. No contraction of the
entire topological group is asserted. Those existence
requirements remain part of C.12.
\end{proof}
